\section{理想层的爆破}\label{sec:ag-2-2}
% 原PDF p23.
\begin{numpara}[Rees 代数]\label{item:ag-2-2-1}
设 $A$ 是环，$I$ 是理想. 与 $(A,I)$ 相伴的\newterm{ag-rees-algebra}{Rees 代数}{Rees algebra}是多项式代数 $A[t]$ 的如下分次 $A$-子代数：
\begin{align*}
\operatorname{Rees}^{I}(A):=\bigoplus_{n\in\mathbb N}I^nt^n=A\oplus It\oplus I^2t^2\oplus\cdots\subset A[t].
\end{align*}
也可简记为 $\operatorname{Rees}^{I}(A)=\bigoplus_{n\in\mathbb N}I^n$，但此时须避免将置于``次数 $1$''的元素 $i\in I$ 与置于``次数 $0$''的同一个元素混淆. 代数 $\operatorname{Rees}^{I}(A)$ 由其次数为 $1$ 的元素生成. 因而有分次 $A$-代数的典范满射
\begin{align*}
\operatorname{Sym}_A(I)\longrightarrow\operatorname{Rees}^{I}(A),
\end{align*}
它将元素 $i\in I=\operatorname{Sym}_A^1(I)$ 映到 $it$.
更一般地，设 $X$ 是概形，$\mathcal I$ 是 $X$ 上的拟凝聚理想层. 记 $\mathcal I^n$ 为乘法诱导的态射 $\mathcal I^{\otimes n}\longrightarrow\mathcal O_X$ 的像. 这是 $\mathcal O_X$ 的拟凝聚理想. 层 $\operatorname{Rees}^{\mathcal I}(\mathcal O_X):=\bigoplus_{n\in\mathbb N}\mathcal I^n$ 是拟凝聚分次 $\mathcal O_X$-代数，也是预层 $U\mapsto\operatorname{Rees}^{\mathcal I(U)}(\mathcal O_X(U))$ 的相伴层. 事实上，对每个仿射开集 $U$，有
\begin{align*}
\Gamma(U,\operatorname{Rees}^{\mathcal I}(\mathcal O_X))=\operatorname{Rees}^{\mathcal I(U)}(\mathcal O_X(U)).
\end{align*}
只需在 $X$ 本身仿射时验证这一点，而为此只需验证 Rees 代数的构造与局部化可交换. 更一般地，对每个平坦态射 $A\longrightarrow A'$，记 $I':=IA'$ 为 $I$ 在 $A'$ 中的像所生成的理想，则有同构 $I^n\otimes_AA'\xrightarrow\sim(I')^n$，由此得到分次 $A'$-代数的同构
\begin{align*}
\operatorname{Rees}^{I}(A)\otimes_AA'\xrightarrow\sim\operatorname{Rees}^{I'}(A').
\end{align*}
由于对称代数的构造与任意换环可交换，还可得到拟凝聚 $\mathcal O_X$-代数的满态射
\begin{align*}
\operatorname{Sym}_{\mathcal O_X}(\mathcal I)\longrightarrow\operatorname{Rees}^{\mathcal I}(\mathcal O_X).
\end{align*}
\end{numpara}
\begin{definition}\label{item:ag-2-2-2}
称 $X$-概形 $\widetilde X_{\mathcal I}:=\operatorname{Proj}(\operatorname{Rees}^{\mathcal I}(\mathcal O_X))\xrightarrow{\pi_{\mathcal I}}X$ 为 $X$ 在拟凝聚理想层 $\mathcal I$ 处的\newterm{ag-blow-up}{爆破}{blow-up}.
\end{definition}
若 $Y$ 是由 $\mathcal I$ 定义的 $X$ 的闭子概形，也记作 $\widetilde X_Y\xrightarrow{\pi_Y}X$，称之为 $X$ 沿 $Y$ 的爆破. 态射 $\pi_Y$ 紧合，但通常不平坦. 注意，若 $A$ 是整环（分别为既约环），则 $A[t]$ 也为整环（分别为既约环），从而无论理想 $I$ 如何，$\operatorname{Rees}^{I}(A)$ 都仍为整环（分别为既约环）. 更一般地，由此得到：若 $X$ 既约（分别为整概形），则 $\widetilde X_{\mathcal I}$ 也既约（分别为整概形）.
此外，还有以下极端情形. 若 $\mathcal I$ 是单位理想（因而定义空子概形），则 $\pi_{\mathcal I}$ 是同构. 若 $\mathcal I$ 幂零，则爆破 $\widetilde X_{\mathcal I}$ 是空概形.
% 原PDF p24--p26上部.
\begin{numpara}[仿射情形的局部描述]\label{item:ag-2-2-3}
假设 $X=\operatorname{Spec}(A)$ 仿射，$\mathcal I$ 与 $A$ 中由元素 $(f_1,\ldots,f_n)$ 生成的理想 $I$ 相伴. 次数为 $1$ 的元素 $f_it\in It=\operatorname{Rees}^{I}_1(A)$ 定义仿射覆盖 $\widetilde X_{\mathcal I}=\bigcup_{i=1}^nD(f_it)$，其中
\begin{align*}
D(f_it)=\operatorname{Spec}(\operatorname{Rees}^{I}(A)/(f_it-1)).
\end{align*}
为更明确地描述此代数，注意包含映射 $\operatorname{Rees}^{I}(A)\subset A[t]$ 诱导从 $\operatorname{Rees}^{I}(A)/(f_it-1)$ 到局部化 $A[t]/(f_it-1)=A_{f_i}$ 的单射（容易验证）. 此外，该单射的像是由元素 $f_jt=f_jf_i^{-1}$ 生成的 $A_{f_i}$ 的 $A$-子代数 $A\langle f_1/f_i,\ldots,f_n/f_i\rangle$. 然而，一般而言，研究并实际计算此代数并不容易.
\end{numpara}
\begin{numpara}[正则理想的爆破]\label{item:ag-2-2-4}
沿用上述设定，考虑将 $T_i$ 映到 $f_it$ 的满射
\begin{align}
\operatorname{Sym}^{\bullet}(A^n)=A[T_1,\ldots,T_n]\twoheadrightarrow\operatorname{Rees}^{I}(A).\tag{2.2.4.1}\label{eq:ag-2-2-4-1}
\end{align}
它诱导闭浸入 $\widetilde X_{\mathcal I}\hookrightarrow\mathbb P^{n-1}\times X$，而由 $D(f_it)$ 组成的覆盖正是 $\mathbb P^n$ 的典范覆盖 $D(T_i)$ 的逆像. 上述满射的核由满足 $P(f_1,\ldots,f_n)=0$（在 $A$ 中）的所有齐次多项式 $P$ 生成. 特别地，它包含所有的 $f_iT_j-f_jT_i$. 我们将证明，当序列 $(f_1,\ldots,f_n)$ 正则时，这些元素生成该核. 先回忆这个概念的定义.
\end{numpara}
\begin{definition*}
若对每个 $i=1,\ldots,n$，$f_i$ 在 $A/(f_1,\ldots,f_{i-1})$ 中的像都是正则元素（即非零因子），则称交换环 $A$ 中的序列 $(f_1,\ldots,f_n)$ 为\newterm{ag-regular-sequence}{正则序列}{regular sequence}.
\end{definition*}
\begin{proposition*}
设 $(f_1,\ldots,f_n)$ 是交换环 $A$ 的正则序列，$I$ 是其生成的理想. 则有：
\begin{enumerate}[label=\roman*)]
\item 分次 $A$-代数态射 $\operatorname{Sym}^{\bullet}(A^n)=A[T_1,\ldots,T_n]\xrightarrow{T_i\mapsto f_it}\operatorname{Rees}^{I}(A)$ 的核由所有的 $f_iT_j-f_jT_i$ 生成.
\item $A$-模态射
\begin{align*}
A^n=\bigoplus_{i=1}^nAT_i\xrightarrow{T_i\mapsto f_i}I
\end{align*}
的核由所有的 $f_iT_j-f_jT_i$ 生成，且典范态射 $\operatorname{Sym}_A^{\bullet}(I)\longrightarrow\operatorname{Rees}^{I}(A)$ 是同构.
\end{enumerate}
\end{proposition*}
后面将看到，若 $A$ 是 Noether 局部环，则满足条件 i) 或 ii) 的每个序列都正则. 此时，正则性这一概念与序列的次序无关.
\begin{proof}
先证明这两个性质等价. 这来自如下的一般事实：设 $\varphi:M\twoheadrightarrow N$ 是 $A$-模满态射，则 $\operatorname{Sym}(\varphi):\operatorname{Sym}_A^{\bullet}(M)\twoheadrightarrow\operatorname{Sym}_A^{\bullet}(N)$ 的核由 $\operatorname{Sym}_A^1(M)=M$ 的子模 $\operatorname{Ker}(\varphi)$ 生成. 例如，可利用对称代数的万有性质验证这一点. 将它应用于态射 $\varphi:A^n=\bigoplus_{i=1}^nAT_i\xrightarrow{T_i\mapsto f_i}I$ 即可. 此外，注意，若 $\operatorname{Ker}(\varphi)$ 确由 $f_jT_i-f_iT_j$ 生成，则模 $I$ 约化后（即与 $A/I$ 作张量积后），$\varphi$ 诱导同构 $(A/I)^n\xrightarrow\sim I/I^2$.
现在对 $n$ 作归纳，证明 i). 记 $J:=(f_1,\ldots,f_{n-1})$. 有分解
\begin{align}
A[T_1,\ldots,T_n]\longrightarrow\operatorname{Rees}^{J}(A)[T_n]\xrightarrow{\psi}\operatorname{Rees}^{I}(A),\tag{2.2.4.2}\label{eq:ag-2-2-4-2}
\end{align}
其中 $\psi$ 将 $T_n$ 映到 $f_nt$，并延拓 $\operatorname{Rees}^{J}(A)$ 到 $\operatorname{Rees}^{I}(A)$ 的典范包含映射. 根据归纳假设，只需验证 $\operatorname{Ker}(\psi)$ 等于 $\operatorname{Rees}^{J}(A)[T_n]$ 中由形如 $jT_n-f_njt$（$j\in J$）的元素生成的理想 $L$.
注意，$\operatorname{Ker}(\psi)$ 是齐次的. 设 $P(T_n)=\sum_{k=0}^Nj_kt^kT_n^{N-k}$ 是 $\operatorname{Ker}(\psi)$ 中次数为 $N$ 的齐次元素（因而 $j_k\in J^k$）. 正则性假设表明，在 $A/J$ 上乘以 $f_n$ 是单射. 因此，由等式 $\sum_kj_kf_n^{N-k}=0$ 可得 $j_0\in J$. 所以模理想 $L$ 而言，$P(T_n)$ 同余于一个多项式 $\sum_{k=1}^Nj_kt^kT_n^{N-k}$. 现在，归纳假设与等价关系 ii) $\Leftrightarrow$ iii) 保证乘法态射 $\operatorname{Sym}_A^m(J)\longrightarrow J^m$ 是同构. 模 $J$ 后，它约化为同构
\begin{align*}
\operatorname{Sym}_{A/J}^m(J/J^2)\xrightarrow\sim J^m/J^{m+1},
\end{align*}
而上面已见 $J/J^2$ 在 $A/J$ 上自由. 因而在 $J^m/J^{m+1}$ 上乘以 $f_n$ 也为单射. 取 $m=1$，等式 $\sum_{k\geqslant1}j_kf_n^{N-k}=0$ 表明 $j_1\in J^2$，从而模 $L$ 而言，$P(T_n)$ 同余于一个多项式 $\sum_{k=2}^Nj_kt^kT_n^{N-k}$. 归纳可知，模 $L$ 而言，$P(T_n)$ 同余于一个``常数'' $j_Nt^N$. 但 $P(f_n)=0$ 蕴含 $j_N=0$.
\end{proof}
回到分目~\ref{item:ag-2-2-3}的记号. 根据命题，当 $(f_i)$ 构成正则序列时，有同构
\begin{align*}
A\left[\frac{f_1}{f_i},\ldots,\frac{f_n}{f_i}\right]\xrightarrow\sim A\left\langle\frac{f_1}{f_i},\ldots,\frac{f_n}{f_i}\right\rangle,
\end{align*}
其中 $A[f_1/f_i,\ldots,f_n/f_i]$ 是 $A[U_1,\ldots,U_n]/(U_1f_i-f_1,\ldots,U_nf_i-f_n)$ 的简写，而同构将 $U_i$ 映到 $A_{f_i}$ 中的元素 $f_jf_i^{-1}$.
\begin{example*}
具体取 $A=k[x_1,\ldots,x_n]$，$I=(x_1,\ldots,x_n)$ 是仿射空间 $X=\mathbb A_k^n$ 的原点 $o$ 处消失的函数所成的理想. 嵌入 $\widetilde{\mathbb A_k^n}_I\hookrightarrow\mathbb A_k^n\times_k\mathbb P_k^{n-1}$ 的像就是其（仿射及齐次）坐标 $(x_1,\ldots,x_n)[T_1:\cdots:T_n]$ 满足等式 $x_iT_j=x_jT_i$ 的点所成的轨迹. 为日后显式计算爆破，宜将 $\widetilde{\mathbb A_k^n}_I$ 看作 $n$ 个仿射坐标图
\begin{align*}
D(x_it)=\operatorname{Spec}\left(k\left[x_i,\frac{x_1}{x_i},\ldots,\frac{x_n}{x_i}\right]\right)\simeq\mathbb A_k^n
\end{align*}
的并. 这些坐标的好处是，它们直接给出坐标图的粘合规则. 例如，$D(x_it)$ 与 $D(x_jt)$ 的交是 $D(x_it)$ 中坐标 $x_j/x_i$ 可逆的仿射开集，即
\begin{align*}
\operatorname{Spec}\left(k\left[x_i,\frac{x_1}{x_i},\ldots,\left(\frac{x_j}{x_i}\right)^{\!\pm1},\ldots,\frac{x_n}{x_i}\right]\right).
\end{align*}
它与 $D(x_jt)$ 中坐标 $x_i/x_j$ 可逆的仿射开集
\begin{align*}
\operatorname{Spec}\left(k\left[x_j,\frac{x_1}{x_j},\ldots,\left(\frac{x_i}{x_j}\right)^{\!\pm1},\ldots,\frac{x_n}{x_j}\right]\right)
\end{align*}
相等同，映射将 $x_i$ 映到 $x_j\cdot(x_i/x_j)$，将 $x_k/x_i$ 映到 $(x_k/x_j)(x_i/x_j)^{-1}$. 此外，态射 $\pi_I$ 在 $D(x_it)$ 上的限制由 $x_j\mapsto x_i\cdot(x_j/x_i)$ 定义. 注意，该限制在原点 $o$ 上方的纤维是 $D(x_it)$ 中方程为 $x_i=0$ 的超平面. 而且，$\pi_I$ 将此超平面在 $D(x_it)$ 中的补集同构地映到它的像，即 $X=\mathbb A_k^n$ 中超平面 $x_i=0$ 的补集. 粘合后可知，$\pi_I$ 在 $o$ 上方的纤维是 $n-1$ 维射影空间，而 $\pi_I$ 诱导同构 $\pi_I^{-1}(X\setminus\{o\})\xrightarrow\sim X\setminus\{o\}$. 纤维 $\pi_I^{-1}(o)$ 称为``\newterm{ag-exceptional-divisor}{例外除子}{exceptional divisor}''. 下面将看到如何推广这些初等的讨论.
\end{example*}
\begin{numpara}[可逆理想的爆破]\label{item:ag-2-2-5}
设 $\mathcal I$ 是 $X$ 上的理想层，且作为 $\mathcal O_X$-模可逆. 局部考察可见，典范态射 $\operatorname{Sym}_{\mathcal O_X}^{\bullet}(\mathcal I)\longrightarrow\operatorname{Rees}^{\mathcal I}(\mathcal O_X)$ 是同构. 因而爆破 $\widetilde X_Y=\mathbb P(\mathcal I)\xrightarrow{\pi_{\mathcal I}}X$ 是同构.
\end{numpara}
\begin{proposition}\label{item:ag-2-2-6}
若 Noether 概形 $X$ 具有丰沛层 $\mathcal L$，则 $\pi_Y$ 是射影态射.
\end{proposition}
\begin{proof}
此时可找到满态射 $\mathcal O_X^{N+1}\twoheadrightarrow\mathcal I\otimes\mathcal L^{\otimes m}$，或者扭转后，令 $\mathcal L':=\mathcal L^{\otimes-m}$，得到满态射 $\mathcal L'^{N+1}\twoheadrightarrow\mathcal I$，从而得到满态射 $\operatorname{Sym}^{\bullet}(\mathcal L'^{N+1})\twoheadrightarrow\operatorname{Rees}^{\mathcal I}(\mathcal O_X)$. 因此有 $X$-嵌入 $\widetilde X_Y\hookrightarrow\mathbb P(\mathcal L'^{N+1})$. 显然有 $\mathbb P(\mathcal L'^{N+1})\simeq\mathbb P(\mathcal O^{N+1})=\mathbb P_X^N$，其态射将 $\mathcal L'^{N+1}\twoheadrightarrow\mathcal F$（$\mathcal F$ 可逆）映到 $\mathcal O^{N+1}\twoheadrightarrow\mathcal F\otimes\mathcal L'^\vee$.
\end{proof}
% 原PDF p26下部--p27上部.
\begin{numpara}[爆破的函子性]\label{item:ag-2-2-7}
从一个交换图出发：
\begin{align}
\begin{tikzcd}[ampersand replacement=\&,column sep=large]
Y'\arrow[r,hook,"i'"]\arrow[d,"g"']\&X'\arrow[d,"f"]\\
Y\arrow[r,hook,"i"']\&X
\end{tikzcd}\tag{2.2.7.1}\label{eq:ag-2-2-7-1}
\end{align}
记 $\mathcal I'$ 为 $Y'$ 在 $X'$ 中的理想. 此图表示态射 $f^*\mathcal I\longrightarrow\mathcal O_{X'}$ 经由 $\mathcal I'$ 分解，由此得到分次 $\mathcal O_{X'}$-代数的态射
\begin{align}
f^*\operatorname{Rees}^{\mathcal I}(\mathcal O_X)\longrightarrow\operatorname{Rees}^{\mathcal I'}(\mathcal O_{X'}).\tag{2.2.7.2}\label{eq:ag-2-2-7-2}
\end{align}
当图~\eqref{eq:ag-2-2-7-1}为笛卡尔图，即 $Y'=X'\times_XY$ 时，态射 $f^*\mathcal I\longrightarrow\mathcal I'$ 满射. 事实上，有关系 $\mathcal I'=f^{-1}(\mathcal I)\cdot\mathcal O_{X'}$，有时称 $\mathcal I'$ 为\newterm{ag-inverse-image-ideal}{逆像理想}{inverse image ideal}. 不应将它与 $\mathcal O_X$-模层 $\mathcal I$ 的逆像 $f^*(\mathcal I)$ 混淆. 在这种情形，$f^*\mathcal I\longrightarrow\mathcal I'$ 的满射性蕴含态射~\eqref{eq:ag-2-2-7-2}满射. 从而得到闭浸入
\begin{align*}
\widetilde{X'}_{Y'}\hookrightarrow X'\times_X\widetilde X_Y.
\end{align*}
左边称为 $X'$ 在 $\widetilde X_Y$ 中的\newterm{ag-strict-transform}{严格变换}{strict transform}，右边称为 $X'$ 在 $\widetilde X_Y$ 中的\newterm{ag-total-transform}{全变换}{total transform}. 上述态射也可由以下交换图描述：
\begin{align}
\begin{tikzcd}[ampersand replacement=\&,column sep=large]
\widetilde{X'}_{Y'}\arrow[r,"\pi_{Y'}"]\arrow[d,"\widetilde f"']\&X'\arrow[d,"f"]\\
\widetilde X_Y\arrow[r,"\pi_Y"']\&X.
\end{tikzcd}\tag{2.2.7.3}\label{eq:ag-2-2-7-3}
\end{align}
\end{numpara}
\begin{lemma}\label{item:ag-2-2-8}
假设图~\eqref{eq:ag-2-2-7-1}为笛卡尔图. 则：
\begin{enumerate}[label=\roman*)]
\item 若 $f$ 平坦，则上图为笛卡尔图，即 $\widetilde{X'}_{Y'}\xrightarrow\sim X'\times_X\widetilde X_Y$.
\item 若 $f$ 是闭浸入，则 $\widetilde f$ 也是闭浸入.
\end{enumerate}
\end{lemma}
\begin{proof}
i) 来自如下事实：若 $f$ 平坦且图~\eqref{eq:ag-2-2-7-1}为笛卡尔图，则 Rees 代数的态射~\eqref{eq:ag-2-2-7-2}是同构. 对于 ii)，注意 $f_g$ 是如下复合：
\begin{align*}
\widetilde{X'}_{Y'}\hookrightarrow X'\times_X\widetilde X_Y\hookrightarrow X\times_X\widetilde X_Y=\widetilde X_Y.\qedhere
\end{align*}
\end{proof}
\begin{numpara}[应用]\label{item:ag-2-2-9}
引理的 ii) 对计算很有用. 例如，设 $k$ 是域，$X$ 是 $\mathbb A_k^n$ 中经过原点 $o$、维数 $>0$ 的子簇（即整的子 $k$-概形）. 根据引理的 ii)，爆破 $\widetilde X_o$ 可等同于 $\widetilde{\mathbb A_k^n}_o$ 的一个闭子概形. 又已知此子概形为整概形，故其非空开集 $\pi_o^{-1}(X\setminus\{o\})$ 在其中稠密. 因此，$\widetilde X_o$ 可以恢复为 $\pi_o^{-1}(X)\setminus\pi_o^{-1}(\{o\})$ 在 $\widetilde{\mathbb A_k^n}_o$ 中的闭包（用文字说，严格变换就是从全变换中去掉例外除子后所得部分的闭包）. 这里指通常意义下的闭包（赋予既约结构的闭子概形）. 计算时，在 $\widetilde{\mathbb A_k^n}_o$ 的``典范''仿射坐标图中进行. 设 $I\subset k[x_i,x_1/x_i,\ldots,x_n/x_i]$ 为全变换 $\pi_o^{-1}(X)\cap D(x_it)$ 的理想. 则严格变换 $\widetilde X_o\cap D(x_it)$ 的理想 $J$ 是以下复合的核：
\begin{align*}
k\left[x_i,\frac{x_1}{x_i},\ldots,\frac{x_n}{x_i}\right]\longrightarrow k\left[x_i^{\pm1},\frac{x_1}{x_i},\ldots,\frac{x_n}{x_i}\right]/I.
\end{align*}
换言之，
\begin{align*}
k\left[x_i,\frac{x_1}{x_i},\ldots,\frac{x_n}{x_i}\right]/J=\left(k\left[x_i,\frac{x_1}{x_i},\ldots,\frac{x_n}{x_i}\right]/I\right)_{/x_i\text{-挠}}.
\end{align*}
后面将看到，即使 $X$ 不是整概形，这些公式也成立.
\end{numpara}
\begin{example*}
当 $n=2$ 时，取 $X$ 为方程 $ax_1+bx_2=0$ 定义的直线. 则 $X$ 的全变换在 $D(x_1t)$ 中的方程是 $x_1(a+b\,x_2/x_1)=0$，严格变换的方程是 $a+b\,x_2/x_1=0$. 与 $D(x_2t)$ 中的情形粘合后，可见 $\pi_o^{-1}(X)$ 是例外除子 $\pi_o^{-1}(o)\simeq\mathbb P_k^1$ 与严格变换的并；后者是一条直线，与 $\mathbb P^1$ 相交于齐次坐标为 $[a:b]$ 的点. 因而，$\mathbb A_k^2$ 在 $o$ 处的爆破按方向``分开''经过 $o$ 的直线.
\end{example*}
% 原PDF p28.
\begin{example*}
仍取 $n=2$，令 $X$ 为方程 $x_2^2=x_1^2(x_1-1)$ 定义的平面三次曲线. 它在 $(0,0)$ 处有一个二重点，该点处的切线为 $x_1+x_2=0$ 和 $x_1-x_2=0$. $X$ 的全变换在 $D(x_1t)$ 中的方程为 $x_1^2((x_2/x_1)^2-x_1+1)=0$. 因而同一坐标图中严格变换的方程为 $(x_2/x_1)^2-x_1+1=0$. 这是非奇异曲线的方程，该曲线与例外除子相交于点 $[1:1]$ 和 $[1:-1]$.
\end{example*}
\begin{example*}
当 $n=3$ 时，取 $X$ 为方程 $x_1x_2=x_3^2$ 定义的曲面. 这是一个锥面，在 $(0,0,0)$ 处有孤立奇点. 可验证，严格变换 $\widetilde X_o$ 在 $D(x_1t)$、$D(x_2t)$、$D(x_3t)$ 中分别由方程
\begin{align*}
\frac{x_2}{x_1}=\left(\frac{x_3}{x_1}\right)^2,\qquad
\frac{x_1}{x_2}=\left(\frac{x_3}{x_2}\right)^2,\qquad
\frac{x_1}{x_3}\cdot\frac{x_2}{x_3}=1
\end{align*}
给出. 注意，最后一个坐标图 $\widetilde X_o\cap D(x_3t)$ 包含在前两个的并中. 所以 $\widetilde X_o$ 由仿射平面
\begin{align*}
\operatorname{Spec}(k[x_1,x_3/x_1])\simeq\mathbb A_k^2,\qquad\operatorname{Spec}(k[x_2,x_3/x_2])\simeq\mathbb A_k^2
\end{align*}
沿同构
\begin{align*}
\operatorname{Spec}\left(k\left[x_1,\left(\frac{x_3}{x_1}\right)^{\!\pm1}\right]\right)\xrightarrow\sim
\operatorname{Spec}\left(k\left[x_2,\left(\frac{x_3}{x_2}\right)^{\!\pm1}\right]\right)
\end{align*}
粘合而成；此同构将 $x_3/x_1$ 映到 $(x_3/x_2)^{-1}$，将 $x_1\cdot(x_3/x_1)$ 映到 $x_2\cdot(x_3/x_2)$. 特别地，$\widetilde X_o$ 是非奇异曲面，且 $\widetilde X_o$ 中的纤维 $\pi_o^{-1}(o)$ 同构于 $\mathbb P_k^1$.
\end{example*}
\begin{theorem}\label{item:ag-2-2-10}
设 $X$ 是（分离）概形，$Y$ 是相伴理想为 $\mathcal I$ 的子概形.
\begin{enumerate}[label=\roman*)]
\item 设 $U$ 是 $Y$ 在 $X$ 中的开补集. $\pi_{\mathcal I}$ 在 $\pi_{\mathcal I}^{-1}(U)$ 上的限制是同构 $\pi_{\mathcal I}^{-1}(U)\xrightarrow\sim U$.
\item $\mathcal I$ 在 $\widetilde X_{\mathcal I}$ 中的逆像理想 $\pi_{\mathcal I}^{-1}(\mathcal I)\cdot\mathcal O_{\widetilde X_{\mathcal I}}$ 典范同构于 $\mathcal O_{\widetilde X_{\mathcal I}}(1)$. 特别地，它（作为模）可逆.
\item 对每个使逆像理想 $f^{-1}(\mathcal I)\mathcal O_{X'}$ 可逆的态射 $f:X'\longrightarrow X$，存在唯一的态射 $g:X'\longrightarrow\widetilde X_{\mathcal I}$，使 $f=\pi_{\mathcal I}\circ g$.
\end{enumerate}
\end{theorem}
可以用闭子概形的语言重新表述 ii)，即：闭子概形 $\pi_Y^{-1}(Y):=Y\times_X\widetilde X_Y$ 是 $\widetilde X_Y$ 上的除子. 此外，对每个使 $Y':=f^{-1}(Y)$ 成为 $X'$ 上除子的态射 $X'\longrightarrow X$，存在唯一的态射 $g:X'\longrightarrow\widetilde X_{\mathcal I}$，使 $f=\pi\circ g$.
除子 $\pi^{-1}(Y)$ 称为例外除子.
\begin{proof}
由于浸入 $U\hookrightarrow X$ 是平坦态射，前面的引理表明，$\pi_{\mathcal I}^{-1}(U)\longrightarrow U$ 可等同于理想 $\mathcal I\mathcal O_U$ 的爆破. 此理想是单位理想，因而可逆，由例~\ref{item:ag-2-2-7}可知这是同构.
对于 ii)，只需局部考察，因而可假设 $X=\operatorname{Spec}(A)$ 仿射. 此时，$I$ 的逆像 $\mathcal O_{\widetilde X_{\mathcal I}}$-理想由分次理想 $I\operatorname{Rees}^{I}(A)=I\bigl(\bigoplus_{n\in\mathbb N}I^n\bigr)$ 给出. 包含映射 $I\operatorname{Rees}^{I}(A)\hookrightarrow\operatorname{Rees}^{I}(A)(1)$ 诱导相伴层的同构，故有 $\pi_{\mathcal I}^{-1}(\mathcal I)\cdot\mathcal O_{\widetilde X_{\mathcal I}}\xrightarrow\sim\mathcal O_{\widetilde X_{\mathcal I}}(1)$.
对于 iii)，记 $\mathcal I':=f^{-1}(\mathcal I)\mathcal O_{X'}$，并令 $Y'=Y\times_XX'$ 为相伴子概形. 在图~\eqref{eq:ag-2-2-7-3}中，态射 $\pi_{Y'}=\pi_{\mathcal I'}$ 是同构，因为 $\mathcal I'$ 可逆. 这就得到所述分解. 还须证明唯一性. 设 $U'$ 为开集 $X'\setminus Y'$. 由 $U'=f^{-1}(U)$ 及 i)，每个使 $f=\pi\circ g$ 成立的 $g$ 在 $U'$ 上都已确定.
因此，只需证明：两个源为 $X'$、目标分离且在 $U'$ 上重合的态射 $g_1,g_2$ 必然相等. 换言之，须说明 $U'$ 概形论稠（参见定理~\ref{item:ag-1-2-4}的证明）. 而 $U'$ 是除子的补集，故概形论稠. [详细地说：先用仿射概形覆盖目标，再覆盖源，便归结到仿射情形. 理想 $\mathcal I'$ 可逆，故可假设它在所考虑的仿射开集 $V'=\operatorname{Spec}(A')$ 上的限制 $I'$ 是主理想 $I'=(f')$，其中 $f$ 正则. 此时 $U'\cap V'=U'_f$ 仍仿射，其环为 $A_f$. 而态射 $A\longrightarrow A_f$ 为单射，因为 $f$ 正则.]
\end{proof}
% 原PDF p29.
\begin{example}\label{item:ag-2-2-11}
假设 $X=\operatorname{Spec}(A)$ 且 $I=(f_1,\ldots,f_n)$. 用仿射开集 $D(f_it)$ 覆盖 $\widetilde X_{\mathcal I}$. 根据定理，例外除子 $E$ 的理想 $\mathcal I'=\pi_{\mathcal I}^{-1}(\mathcal I)\mathcal O_{\widetilde X_{\mathcal I}}$ 局部为主理想，由一个正则截面生成. 更确切地说，开集 $D(f_it)$ 恰好由使截面 $\pi_{\mathcal I}^*f_i$ 正则且生成此理想的点组成.\unskip\nobreak\hfill\ensuremath{\diamond}
\renewcommand{\envenddiamond}{}% final marker on the text line
\end{example}
\begin{numpara}[应用：严格变换是从全变换中去掉例外除子后所得部分的概形论闭包]\label{item:ag-2-2-12}
这推广了分目~\ref{item:ag-2-2-9}的讨论，并说明其中给出的计算即使不假设 $X$ 为整概形也仍然有效.
先回忆概形论闭包的概念. 更一般地，给定 Noether 概形的态射 $Y\xrightarrow f X$，称由理想层 $\operatorname{Ker}(\mathcal O_X\longrightarrow f_*\mathcal O_Y)$ 定义的 $X$ 的闭子概形为 $f$ 的\newterm{ag-scheme-theoretic-image}{概形论像}{scheme-theoretic image}. 它的支集是 $f$ 的点集像的拓扑闭包. 当 $Y$ 是 $X$ 的子概形（因而局部闭，或者说 $f$ 是浸入）时，也称 $f$ 的概形论像为 $Y$ 的\newterm{ag-scheme-theoretic-closure}{概形论闭包}{scheme-theoretic closure}. 例如，设 $S$ 是离散赋值谱（trait），即离散赋值环 $R$ 的谱，$X$ 是 $S$-概形. 则 $X$ 的泛纤维（即 $S$ 的泛点在 $X$ 中的原像）的概形论闭包，是由 $\mathcal O_X$ 的 $R$-挠理想定义的 $X$ 的闭子概形.
现在，设 $X$ 是 Noether 概形，$Y,X'$ 是 $X$ 的两个闭子概形，$Y':=Y\times_XX'$. 由引理~\ref{item:ag-2-2-8}的 ii)，$\widetilde{X'}_{Y'}$ 在 $\widetilde X_Y$ 中闭. 而且，根据上述定理，$\pi_{Y'}^{-1}(Y')$ 是 $\widetilde{X'}_{Y'}$ 中的除子，故其开补集在 $\widetilde{X'}_{Y'}$ 中概形论稠. 这表明 $\widetilde{X'}_{Y'}$ 可等同于 $\pi_Y^{-1}(X')\setminus\pi_Y^{-1}(Y)$ 的概形论闭包.
\end{numpara}
\begin{corollary}\label{item:ag-2-2-13}
设 $\mathcal I$ 和 $\mathcal J$ 是 $X$ 上的两个理想层. 则有分解
\begin{align*}
\pi_{\mathcal I\mathcal J}:\widetilde X_{\mathcal I\mathcal J}\xrightarrow{\pi}\widetilde X_{\mathcal I}\xrightarrow{\pi_{\mathcal I}}X,
\end{align*}
其中 $\pi$ 可等同于 $\widetilde X_{\mathcal I}$ 中理想 $\pi_{\mathcal I}^{-1}(\mathcal J)\mathcal O_{\widetilde X_{\mathcal I}}$ 的爆破.
\end{corollary}
注意两个特例：
\begin{itemize}
\item 若 $\mathcal J$ 可逆，则 $\pi$ 是同构.
\item 对 $d\in\mathbb N$，爆破 $\widetilde X_{\mathcal I^d}$ 与 $\widetilde X_{\mathcal I}$ 在 $X$ 上同构.
\end{itemize}
\begin{proof}
先验证，两个理想的乘积可逆，当且仅当两个理想各自可逆. 只须证明必要性. 问题是局部的，设 $A$ 为局部环，$I,J$ 为两个理想，且乘积为主理想 $IJ=fA$，其中 $f$ 正则. 可写成 $f=\sum_ki_kj_k$，且对每个 $k$ 存在 $a_k$ 使 $i_kj_k=a_kf$. 因为 $f$ 正则，有等式 $\sum_ka_k=1$；又因 $A$ 局部，某个 $a_k$，不妨记为 $a_0$，是单位. 将 $i_0$ 换成 $i_0a_0^{-1}$，便有 $i_0j_0=f$. 特别地，$i_0,j_0$ 都正则. 取 $I$ 中的元素 $i$，写成 $ij_0=af$. 约去正则元素 $j_0$，得 $i=i_0a$，故 $I=i_0A$ 可逆. 同理，$J$ 可逆.
现在，命题所述分解的存在性来自爆破 $\pi_{\mathcal I}$ 的万有性质. 同样由此性质可得 $\widetilde X_{\mathcal I}$-态射
\begin{align*}
f:\widetilde X_{\mathcal I\mathcal J}\longrightarrow\widetilde{(\widetilde X_{\mathcal I})}_{\mathcal J},
\end{align*}
其中右边表示 $\widetilde X_{\mathcal I}$ 在理想 $\pi_{\mathcal I}^{-1}(\mathcal J)\mathcal O_{\widetilde X_{\mathcal I}}$ 处的爆破. 由下面的引理，$\mathcal I$ 在 $\widetilde{(\widetilde X_{\mathcal I})}_{\mathcal J}$ 上的逆像理想可逆. 事实上，它局部为主理想（因为 $\mathcal I$ 如此），且它在例外除子的补集上的限制可逆. 因而得到 $\widetilde X_{\mathcal I}$-态射
\begin{align*}
g:\widetilde{(\widetilde X_{\mathcal I})}_{\mathcal J}\longrightarrow\widetilde X_{\mathcal I\mathcal J}.
\end{align*}
复合 $g\circ f$ 与 $f\circ g$ 在概形论稠的开集上与恒等态射重合，因而就是恒等态射，这证明了第二个断言.
对于第一个特例，应用下面的引理，以与上面相同的论证可知理想 $\pi_{\mathcal I}^{-1}(\mathcal J)\mathcal O_{\widetilde X_{\mathcal I}}$ 可逆. 第二个特例可对 $d$ 作归纳得到，但更简单的证明是注意到 $\operatorname{Rees}^{\mathcal I^d}(\mathcal O_X)=(\operatorname{Rees}^{\mathcal I}(\mathcal O_X))^{(d)}$，参见分目~\ref{item:ag-2-1-1}.
\end{proof}
\begin{lemma*}
设 $X$ 是概形，$\mathcal I$ 是 $X$ 上局部为主理想的理想. 若 $\mathcal I$ 在一个除子的开补集 $U$ 上的限制可逆（更一般地，当 $X$ 为 Noether 概形时，可取概形论稠的开集 $U$），则 $\mathcal I$ 可逆.
\end{lemma*}
\begin{proof}
问题是局部的，故可假设 $X=\operatorname{Spec}(A)$，其中 $A$ 局部. 理想 $I$ 由某个元素 $f$ 生成，须说明此元素正则. 设 $h\in A$ 满足 $fh=0$. 根据假设，$U$ 具有形式 $\operatorname{Spec}(A[g^{-1}])$，其中 $g$ 正则. （在 Noether 情形，每个概形论稠的开集 $U$ 都包含形如 $\operatorname{Spec}(A[g^{-1}])$ 的概形论稠开集，参见 EGA IV.20.2.12.）由 $I$ 在 $U$ 上可逆，还可在进一步缩小 $U$ 后，假设 $IA[g^{-1}]$ 由正则元素 $f'$ 生成. 因为 $f'\in fA[g^{-1}]$，有 $f'h=0$. 由正则性，$h$ 在 $A[g^{-1}]$ 中为零. 而 $g$ 也正则，所以 $h$ 在 $A$ 中为零.
\end{proof}
注意：$\widetilde X_{\mathcal I^d}=\widetilde X_{\mathcal I}$ 并不意味着 $I$ 的爆破与 $\sqrt I$ 的爆破相同. 例如，在 $\mathbb A_k^2=\operatorname{Spec}(k[x,y])$ 中爆破理想 $(x^2,y)$. 序列 $(x^2,y)$ 正则，所以可以应用分目~\ref{item:ag-2-2-4}. 可见 $\widetilde{\mathbb A_k^2}_I$ 由两个仿射坐标图覆盖：
\begin{align*}
D(x^2t)&=\operatorname{Spec}\left(k\left[x,y,\frac y{x^2}\right]\right)=\operatorname{Spec}\left(k\left[x,\frac y{x^2}\right]\right)\simeq\mathbb A_k^2,\\
D(yt)&=\operatorname{Spec}\left(k\left[x,y,\frac{x^2}y\right]\right)=\operatorname{Spec}(k[x,y,v]/(x^2-yv)),
\end{align*}
后者是一个锥面，在 $(0,0,0)$ 处有二次奇点. 特别地，$\widetilde{\mathbb A_k^2}_I$ 是奇异的. 还应注意，例外除子在 $D(x^2t)$ 中由理想 $(x^2)$ 给出，在 $D(yt)$ 中由理想 $(y)$ 给出，是``加厚''的 $\mathbb P_k^1$，同构于 $\mathbb P^1\times\operatorname{Spec}(k[x]/(x^2))$. 相反，$o$ 的逆像与 $(x,y)$ 的爆破一样，是 $\mathbb P_k^1$.
上述推论有如下逆命题.
% 原PDF p31--p32（止于2.3之前）.
\begin{proposition}\label{item:ag-2-2-14}
有限型理想的两次爆破的复合仍是爆破.
\end{proposition}
\begin{example*}
设 $X=\mathbb A_k^2$，$X'$ 是 $X$ 在理想 $(x,y)$ 处的爆破，$X''$ 是 $X'$ 在仿射坐标图 $D(xt)$ 中由理想 $(x,y/x)$ 给出的点处的爆破. 这也是 $X'$ 在同一仿射坐标图的理想 $(x)(x,y/x)=(x^2,y)$ 处的爆破. 因而，$X''$ 可等同于 $(x,y)(x^2,y)=(x^3,xy,y^2)$ 的爆破.
\end{example*}
\begin{proof}
（证明梗概，参见 Raynaud--Gruson，引理 5.1.4.）考虑复合
\begin{align*}
\widetilde{(\widetilde X_{\mathcal I})}_{\mathcal J}\xrightarrow{\widetilde\pi_{\mathcal J}}\widetilde X_{\mathcal I}\xrightarrow{\pi_{\mathcal I}}X,
\end{align*}
其中 $\mathcal J$ 是 $\mathcal O_{\widetilde X_{\mathcal I}}$ 的有限型拟凝聚理想. 根据前面的推论，只需找到整数 $n$，使 $\mathcal O_{\widetilde X_{\mathcal I}}$ 的理想 $\pi_{\mathcal I}^{-1}(\mathcal I^n)\mathcal J$ 具有形式 $\pi_{\mathcal I}^{-1}(\mathcal K)\mathcal O_{\widetilde X_{\mathcal I}}$，其中 $\mathcal K$ 是 $\mathcal O_X$ 的理想.
为简便起见，假设 $X=\operatorname{Spec}(A)$ 仿射. 此时由上面的定理，$\pi_{\mathcal I}^{-1}(\mathcal I)\mathcal O_{\widetilde X_{\mathcal I}}$ 丰沛\footnote{对一般的 $X$，此理想在 EGA II.4.6 的意义下``相对丰沛''，利用 EGA II.4.6.8 也能同样处理.}. 因而可找到 $n,N$ 及满态射 $\mathcal O_{\widetilde X_{\mathcal I}}^N\twoheadrightarrow\pi_{\mathcal I}^{-1}(\mathcal I^n)\mathcal J$. 由伴随关系，得到
\begin{align*}
\mathcal O_X^N\longrightarrow\pi_{\mathcal I,*}(\pi_{\mathcal I}^{-1}(\mathcal I^n)\mathcal J)\subset\pi_{\mathcal I,*}(\mathcal O_{\widetilde X_{\mathcal I}}(n)),
\end{align*}
即 $\Gamma(\widetilde X_{\mathcal I},\mathcal O_{\widetilde X_{\mathcal I}}(n))$ 中的 $N$ 个元素. 问题在于，典范态射 $I^n\longrightarrow\Gamma(\widetilde X_{\mathcal I},\mathcal O_{\widetilde X_{\mathcal I}}(n))$ 一般并不满射. 然而，如定理~\ref{item:ag-2-1-11}之后的注所述，当 $n$ 充分大时它满射. 于是，只需取 $K$ 为包含在 $I^n$ 中、由 $I^n$ 中像为上述 $N$ 个截面的 $N$ 个元素生成的 $A$ 的理想.
\end{proof}
上述结果和例子表明，爆破的行为多种多样（可以简化奇点，也可能使奇点更严重，等等）. 若只考虑整概形，下述定理表明，爆破理想是一个非常（甚至过于）一般的方法.
\begin{theorem}\label{item:ag-2-2-15}
设 $f:Z\longrightarrow X$ 是域 $k$ 上拟射影簇之间的射影双有理态射（更一般地，可假设 $Z,X$ 都是 Noether 整概形，且 $X$ 具有丰沛层）. 则 $f$ 同构于 $\mathcal O_X$ 的某个理想的爆破.
\end{theorem}
证明之前，先注意，对整概形 $X$，也可以爆破分式理想（即 $\mathcal M_X$ 的有限型 $\mathcal O_X$-子模）. 事实上，对这样的理想 $\mathcal I$，分次 $\mathcal O_X$-代数 $\operatorname{Rees}^{\mathcal I}(\mathcal O_X)=\bigoplus\mathcal I^n$ 仍拟凝聚. 其 $\operatorname{Proj}$ 记为 $\widetilde X_{\mathcal I}$；它是整概形，投影 $\pi_{\mathcal I}$ 双有理，逆像 $f^{-1}(\mathcal I)\mathcal O_{\widetilde X_{\mathcal I}}$ 是 $\mathcal M_{\widetilde X_{\mathcal I}}$ 的可逆分式理想，而 $\widetilde X_{\mathcal I}$ 对这些性质是万有的. 若 $X$ 还具有丰沛层，则可找到 $\mathcal O_X$ 的可逆理想 $\mathcal K$，使 $\mathcal K\mathcal I\subset\mathcal O_X$（取一个到 $\mathcal I$ 的分母层 $\mathcal J$ 有非零态射的 $\mathcal K$）. 此时，爆破 $\pi_{\mathcal I}:\widetilde X_{\mathcal I}\longrightarrow X$ 同构于``通常的''爆破 $\pi_{\mathcal K\mathcal I}:\widetilde X_{\mathcal K\mathcal I}\longrightarrow X$，所以这种爆破可逆理想的方法并不比爆破通常理想更一般.
\begin{proof}
\begin{enumerate}[label=\arabic*-]
\item 先证明，$Z$ 上存在可逆层 $\mathcal L$，使 $Z\simeq\operatorname{Proj}(S_{\mathcal L,\bullet})$，其中 $S_{\mathcal L,\bullet}=\mathcal O_X\oplus\bigoplus_{n>0}f_*(\mathcal L^{\otimes n})$ 由次数为 $1$ 的部分生成.
为此，先取相对于 $X$ 极丰沛的 $\mathcal L$. 已知每个 $f_*(\mathcal L^{\otimes n})$ 在 $X$ 上凝聚（定理~\ref{item:ag-2-1-11}），但尚不知 $S_{\mathcal L,\bullet}$ 是否由次数为 $1$ 的部分生成. 先假设 $X=\operatorname{Spec}(A)$ 仿射. 则 $Z$ 形如 $\operatorname{Proj}(S_\bullet)$，其中 $S_\bullet$ 由次数为 $1$ 的部分生成，且 $\mathcal O(1)=\mathcal L$. 此时，代数 $S_{\mathcal L,\bullet}$ 可等同于 $S'=\Gamma_\bullet(\mathcal O)=\bigoplus_n\Gamma(Z,\mathcal O(n))$. 在定理~\ref{item:ag-2-1-11}之后的注中，已经看到 $S'$ 与 $S$ 在充分高的次数上一致，所以存在整数 $d$，使 $\operatorname{Proj}(S'^{(d)})\xrightarrow\sim\operatorname{Proj}(S)=Z$. 而 $S'^{(d)}$ 正是 $S_{\mathcal L^{\otimes d},\bullet}$. 回到拟射影情形，可用有限多个仿射开集覆盖 $A$，由此得到一个整数 $d$，使 $\mathcal L^{\otimes d}$ 满足要求.
\item 因为 $Z$ 为整概形，$\mathcal L$ 可嵌入 $\mathcal M_Z$，从而有嵌入 $f_*\mathcal L\hookrightarrow f_*\mathcal M_Z$. 又因 $f$ 双有理，有 $f_*\mathcal M_Z=\mathcal M_X$. 所以 $f_*\mathcal L$ 同构于一个分式理想层（它先验未必可逆）.
\item $Z\xrightarrow f X$ 同构于 $f_*(\mathcal L)$ 的爆破.
须将 $S_{\mathcal L,\bullet}$ 与 $\operatorname{Rees}^{f_*(\mathcal L)}(\mathcal O_X)$ 等同，而为此须将 $f_*(\mathcal L^{\otimes n})$ 与分式理想 $f_*(\mathcal L)$ 的 $n$ 次幂 $(f_*(\mathcal L))^n$ 等同. 首先，可将 $\mathcal L^{\otimes n}$ 与分式理想 $\mathcal L$ 的 $n$ 次幂 $\mathcal L^n$ 等同. 于是有包含关系 $f_*(\mathcal L)^n\subset f_*(\mathcal L^n)$. 另一方面，$S_{\mathcal L,\bullet}$ 由次数为 $1$ 的部分生成，所以有满态射 $f_*(\mathcal L)^{\otimes n}\twoheadrightarrow f_*(\mathcal L^{\otimes n})$. 因而有分式理想的等式 $f_*(\mathcal L)^n=f_*(\mathcal L^n)$.\qedhere
\end{enumerate}
\end{proof}
